% secV8_indicators.tex -- Section 9. CONTAINER (27 Oct 2026): text to be written by the author.
% Word budget: 1,450 words (JPART count: running text, footnotes and citations as printed).
% Outline, numbers and sources: referee_EE_2026-10/ESQUELETO.md, section 9. The archived text
% (_arquivo_texto_IA_2026-10-27/) is not a source.

\section{Making a childcare target verifiable}\label{sec:indicators}

The measurement problem is no longer confined to S\~ao Paulo. Brazil's new National Education Plan, enacted by Law 15,388/2026 \citep{brasil2026pne}, requires provision for 100\% of \emph{demanda manifesta} (expressed demand) for childcare and, nationally, coverage of at least 60\% of children aged three or younger. The law does not define expressed demand. Article~28 assigns the National Institute for Educational Studies and Research (Inep) 180 days from the law's publication on April 15, 2026 to establish indicators and make their conceptual bases public \citep{brasil2026pne}. In September 2026, the Ministry of Education established a working group to develop the indicators and projections, initially for three months, extendable by another three (Ordinance 701/2026) \citep{mec2026portaria701}. The first monitoring publication is scheduled after 18 months (Article~30) \citep{brasil2026pne}. We had not located the new Inep indicators by October 26, 2026, without establishing whether the process had been completed.

The national target reproduces Goal~5's measurement choices: the definition of a request and its resolution, observation dates, enrolment source and network, population series, and the relation between demand and coverage requirements. A percentage target alone does not determine an indicator.

\subsection{From a dated stock to the time taken to obtain a place}

Unserved expressed demand concerns a family that requests care but receives no place when it is needed. A waiting-list stock on one date is an imperfect proxy. December follows the main referral round, whereas September includes recent requests. Neither count distinguishes a request for the first available place from a decision to await a particular centre. As Section~\ref{sec:service} shows, the age at first registration also changed. An identical list count can thus represent different waiting experiences.

A more direct measure follows requests from their submission date. Let $\mathcal{R}_s$ be the cohort of eligible first requests registered during period $s$, and let $T$ be a publicly specified interval in days. We propose reporting
\begin{equation}
F_s(T)=\Pr\!\left(\text{enrolment within }T\text{ days of first request}\mid i\in\mathcal{R}_s\right),
\label{eq:timely_enrollment}
\end{equation}
with enrolment as the event of interest and ageing out and definitive exits as competing events. The unserved share is $1-F_s(T)$; reporting should distinguish waiting from exit without enrolment. Eligibility, refusals, and follow-up rules must be stated; cohorts without $T$ days of potential follow-up should not be treated as fully observed. The unit is the child, not an administrative application line. Referral alone is not service received: the allocating agency can issue a referral without the child enrolling. Where possible, reported waiting time should also be related to the family's desired starting date, which is absent from the present records.

A jurisdiction may instead retain a stock-based performance measure. In that case, publishing its annual average alongside quarterly or monthly observations would reduce sensitivity to the choice of a single snapshot, although it would not resolve the conceptual problem. Restricting the stock to requests older than 30 days, as in S\~ao Paulo's municipal Goal~77 \citep{sme_metas2025}, still leaves the reference date consequential.

The contrast is measurable in the pre-2019 records. Between 2014 and 2017, the published stock measured in March, June, September, and December fell by 20.5\%, 18.9\%, 17.5\%, and 53.2\%, respectively (Table~\ref{tab:stab}). The ratio of the December to September percentage declines is 3.03. For the mean of the four published counts over the preceding 12 months, the corresponding declines are 20.8\%, 20.8\%, 21.5\%, and 25.4\%, giving a ratio of 1.18. The same comparison using windows of first requests, registration age, or referral within 365 days produces ratios between 0.75 and 1.11. First requests decline by 5.4--8.3\% across the four windows; median age at registration declines by 3.8--5.2 months; cumulative referral incidence rises by 0.130--0.191. The single-date stock is substantially more sensitive to the chosen month. The comparison partly reflects the difference between a one-day stock and 12-month flows. The request-based measure cannot be calculated for 2019--2022, when compliance classifications diverge, because the files end in 2018.

\subsection{Four complementary indicators}

A monitoring system should retain both access and timeliness. First, report the share of a birth cohort enrolled in childcare before ages $a=12,24,$ and $36$ months, using a stated birth-cohort denominator and enrolment universe. This is a population-oriented access measure, distinct from the current Inep population indicator based on household survey data \citep{inep2024quintociclo}. As an illustration, among the linked cohorts in our records, enrolment before age 24 rises from 31.9\% for children born in 2012 to 53.2\% for those born in 2016 (Table~\ref{tab:cohortcov}). These register-based cohort ratios are not automatically corrected to the Demographic Census population.

Second, report $F_s(T)$ for fixed values such as 180 and 365 days, distinguishing enrolment from referral and publishing the competing-risk and censoring rules. Third, publish the distribution of age at first request, including the median and shares applying before six, 12, and 24 months. A timeliness rate can improve simply because families begin applying earlier; the age distribution makes that compositional change visible. Fourth, publish a monthly reconciliation of the outstanding stock:
\begin{equation}
L_m=L_{m-1}+I_m-O_m+A_m,
\label{eq:stock_identity}
\end{equation}
where $I_m$ and $O_m$ are, respectively, entries into and exits from the \emph{defined stock}, and $A_m$ records adjustments due to corrections, reclassification, and definition changes. Exits should separately identify enrolments, referrals without enrolment, withdrawals, refusals, ageing out, and other administrative dispositions. Entries into the stock need not coincide with first requests, so both series should be released. An unexplained residual should be reported rather than absorbed into an exit category.

Coverage concerns children, $F_s(T)$ concerns timeliness, registration age concerns applicant composition, and the identity permits stock reconciliation. None measures all childcare needs.

\subsection{Publication rules and institutional accountability}\label{sec:policy}

These results imply a publication protocol. Each performance measure should identify its numerator, denominator, administrative source, event definition, and applicable network and age group. The 2023 municipal review published no formula for its Goal~5 indicator \citep{saopaulo2023conferencia}; a ratio reconstructed from December enrolment and waiting-list counts matches six of eight reported values but does not test the binding legal requirement (Section~\ref{sec:goal}). The formula is therefore necessary for verification.

Counts should be issued with both observation and extraction dates and with observations on either side of the main allocation round. The published list, the list excluding centre-specific preferences, the subset subject to document checks, and the subset waiting beyond a duration threshold should be separately reported at the child level, with changes in definitions documented and overlapping series released at each transition. In the earlier years, the stock reconstructed from individual files was 4--13\% lower than the published count \citep{sme_demanda_escolar}, and the two available list definitions later yielded different compliance findings. Monthly entries and dispositions would make it possible to locate such differences rather than attribute them to an undifferentiated category of ``demand.''

Enrolment should be available by network, age, and reference date, with a reconciliation to the School Census. In 2021, the municipal December enrolment publication counted 385,100 children \citep{sme_demanda_escolar}, more than the 377,862 reported across \emph{all} networks in the May School Census \citep{inep_censo_escolar_c}; substituting one source for another changes a compliance classification. Every population denominator should similarly carry the producing institution, release vintage, age definition, and interpolation rule. The law does not designate the demographic denominator. In 2025, for example, the SME referred to approximately 566,000 children using SEADE data \citep{sme2025filazerada}, while the interpolation of SEADE projections available for this exercise yields approximately 504,000 \citep{seade_idade_escolar}. Even without treating either as erroneous, the discrepancy precludes an unqualified comparison of coverage percentages. The annual local demand survey required by Law 14,851/2024 \citep{brasil2024lei14851} could supply a common reporting schedule, provided its results are linked transparently to the register and enrolment records. Periodic audits of application coverage against enrolment records are also necessary, given the discontinuity established in Section~\ref{sec:coverage}.

Reference dates are an explicit design choice in other public services. England counts households on local-authority social-housing registers on March 31 \citep{mhclg_table600}; Norwegian municipal childcare waiting lists use December 15 \citep{ssb13632}; and NHS England reports incomplete referral-to-treatment pathways at the end of each reporting period \citep{nhs_rtt_mar26}. Such conventions make series reproducible, not necessarily immune to timing effects. The NHS referral-to-treatment measure rose from 62.6\% in February to 65.3\% in March 2026 \citep{nhs_rtt_mar26}, against an interim March target of 65\% \citep{nhs_annualreport2026}. This coincidence does not establish an administrative response to the target date, particularly because the NHS indicator differs from a childcare register.

Finally, a demand requirement cannot replace a population requirement. Registering is itself a selected outcome: in the 2015 national survey, only 43.2\% of respondents seeking childcare for a child not enrolled reported taking some action to obtain a place \citep{ibge2017cuidados}. Barriers faced by nonapplicants will not be revealed by a perfectly administered list. The Supreme Federal Court has recognized an individually enforceable right to early-childhood education (Theme 548) \citep{stf2022tema548}, but this does not settle who can compel compliance with an aggregate numerical target. Retaining two measures guards against treating low registration as universal access.

The proposed definitions do not eliminate discretion. Choosing $T$, deciding which applications are eligible, and specifying what qualifies as enrolment are policy decisions. The gain is that these decisions become inspectable. A published indicator defining service as enrolment within 180 days can be checked and contested; an unqualified assertion that 100\% of expressed demand has been served cannot.
